\documentclass[openany,fontset=fandol]{ctexbook}
\usepackage{amsmath,amssymb}
\usepackage[unicode,hidelinks]{hyperref}
\setlength{\emergencystretch}{2em}
\interfootnotelinepenalty=10000
\DeclareMathOperator{\depth}{depth}
\DeclareMathOperator{\Spec}{Spec}
\newcommand{\OO}{\mathcal O}
\newcommand{\FF}{\mathcal F}
\newcommand{\JJ}{\mathcal J}
\raggedbottom
\begin{document}
\title{代数几何基础第五卷预备笔记\\线性系的唯一性与深度}
\author{据法文扫描和英译编辑稿互校续译}
\date{2026年11月22日}
\maketitle
\chapter*{范围与接续}
第79批续译第三篇第3小节，从无固定分支的线性系唯一性命题起，至其最后一项的证明结束。法文物理131--132页（手写138--139），英译物理99--100页。
译校补证明示了原稿省略的非诺特局部论证，并给出最大定义域处的核查理由；与译文区分。
下批从法文物理133页、英译物理100页下部作者反省及命题3.3起。本批尚未并入1158页合订本。
\chapter*{第三篇（续）}
\section*{3\quad 无固定分支线性系的唯一性}

\noindent\textbf{命题3.2。}
设 \(X\to S\) 平坦且局部有限表示，纤维满足 \((S_2)\)。
设 \(D\) 是 \(X/S\) 上由 \(P^\vee\) 参数化的、无固定分支的除子线性系。
记 \(Z\) 为其固定点集。对应的态射 \(f:X\setminus Z\to P\) 决定整个 \(D\)；
甚至 \(f\) 作为“相对于 \(S\) 的伪态射”的等价类也决定 \(D\)。
该等价类的定义域正是 \(X\setminus Z\)。\footnote{法文物理131页命题号附近有多处改笔，英译物理99页不编号。法132及英100在核查末句时都回指“命题2”，下一命题标3.3，故这里按本节编号定为3.2。}

有关“相对于 \(S\) 的伪态射”的概念和基本事实，参见原稿所指的第20.10小节。\footnote{法131手写第20.10小节，取代被划去的原引用；英译脚注77说正式 EGA IV 的对应编号只到20.6。因此保留草稿引用，不把20.10冒充正式出版的 EGA IV 条目。法稿还把“相对有理映射”修改为“相对伪态射”。本段使用的等价关系由紧随其后的共同相对概形稠密开集描述。}
我们要证明：若 \(D'\) 是另一个由 \(P^\vee\) 参数化的无固定分支线性系，
它确定 \(f':X\setminus Z'\to P\)，且 \(f\) 与 \(f'\) 在一个相对于 \(S\) 概形稠密的开集
\[
 U\subset (X\setminus Z)\cap(X\setminus Z')
\]
上一致，那么 \(D=D'\)。

事实上，因 \(P\) 在 \(S\) 上分离，两态射在整个交集上一致，所以可以取
\[
 U=(X\setminus Z)\cap(X\setminus Z')=X\setminus Z'',
 \qquad Z''=Z\cup Z'.
\]
\(Z''\) 在每条纤维中都有余维至少二，且 \(X\) 的纤维满足 \((S_2)\)，所以对每个
\(x\in Z''\)，若 \(s\) 是其在 \(S\) 中的像，则 \(\depth\OO_{X_s,x}\geq2\)。
再利用 \(X\) 在 \(S\) 上平坦且局部有限表示，得到：
\(X\) 上的除子（不必横截各纤维）由它在 \(X\setminus Z''\) 上的限制决定。
把这一结论用于 \(X\times_S P^\vee\to P^\vee\) 及相应闭补集，就得到所需结论。\footnote{法131明确写 \(X-Z''\)，英100先误作 \(Z-Z''\)，后又漏了两个撇号。线性系的两个除子实际在 \(X\times_S P^\vee\) 上，故正文写明基变换。域扩张保持纤维余维，并且纤维在域扩张下仍满足 \((S_2)\)；也可直接用深度沿平坦局部同态不下降。两个态射的相等性由分离态射的闭对角与概形稠密性得到，不是仅凭拓扑稠密。}

设 \(\JJ\) 是定义有效除子的理想。只须证明典范映射
\[
 \JJ\longrightarrow i_*(\JJ|_{X\setminus Z''})
\]
是同构，其中 \(i:X\setminus Z''\hookrightarrow X\) 是典范开浸入。
事实上，由同一个结论也可恢复 \(\OO_X\)；
因此对 \(\JJ|_{X\setminus Z''}\to\OO_X|_{X\setminus Z''}\) 施加 \(i_*\)，也就恢复出包含映射 \(\JJ\to\OO_X\)。
\(\JJ\) 可逆，因而在 \(S\) 上平坦；其纤维限制在 \(Z''\) 上各点的深度至少为二。
于是只须证明下面的引理。\footnote{即使不要求有效 Cartier 除子与纤维横截，其理想作为可逆模的纤维限制仍有定义；这里并不把这个限制当成纤维结构层的子理想。本证明只用 \(\JJ|_{X_s}\) 的深度。法稿物理132页把这一段及引理括起，旁记拟移至20.11；没有把它当作已删正文。}

\medskip\noindent\textbf{引理3.1。}
设 \(X\to S\) 局部有限表示，\(\FF\) 是 \(X\) 上有限表示且在 \(S\) 上平坦的拟凝聚模，\(T\subseteq X\) 为闭子集。
设对每个 \(x\in T\)，记 \(s\) 为 \(x\) 的像，并有
\(\depth_{\OO_{X_s,x}}\FF_{s,x}\geq1\)（分别地，\(\geq2\)）。
则典范同态
\[
 \FF\longrightarrow i_*(\FF|_{X\setminus T}),\qquad
 i:X\setminus T\hookrightarrow X,
\]
为单射（分别地，为同构）。\footnote{两稿陈述一度把深度条件写在 \(x\in X\) 上，但法132的证明明确只用 \(x\in T\)；取整个 \(X\) 会排除通常的纤维泛点，并非本命题所需。开浸入的靶是 \(X\)，不是两底本原式中的 \(T\)；英译也在括号中指出。局部有限表示依照法稿。}

\emph{引理的证明。}
可以把 \(S\) 与 \(X\) 局部取为仿射，并按惯用方法归结到底 \(S\) 诺特的情形。
此时由第6段，对于在 \(s\in S\) 上方的任一点 \(x\)，
\[
 \depth_{\OO_{X,x}}\FF_x\ \geq\ \depth_{\OO_{X_s,x}}\FF_{s,x}.
\]
于是 \(x\in T\) 时，左端深度至少为一（分别地，至少为二）；再由第5段即得结论。
作者备注：若要安排妥当，这条引理应当放在第11段“去掉诺特性假设”的标题下。\footnote{法132证明末行仍有纤维下标残留；逐点不等式推出的是总空间模在闭补集处的深度界，才能应用这里的延拓判据。英译者提示亦可参看 EGA IV 11.3。原稿简写的降阶见下方补证。}

最后验证命题3.2的末一断言：\(X\setminus Z\) 正是 \(f\) 所定义的相对伪态射的定义域。
设 \(V\supset X\setminus Z\) 是它的定义域，带有代表态射 \(V\to P\)。
按命题3.1，这个态射在 \(V/S\) 上给出由 \(P^\vee\) 参数化的无固定点线性系 \(D'\)；
这里逐纤维检查命题3.1的条件：新增点位于 \(Z\)，而 \(Z_s\) 不含任何纤维关联点；
在其余开集上的非退化条件因此也排除了在新增点处出现新的退化。
于是
\[
 D'|_{(X\setminus Z)\times_S P^\vee}
 = D|_{(X\setminus Z)\times_S P^\vee}.
\]
利用已证的唯一性，得到 \(D'=D|_{V\times_S P^\vee}\)。
所以 \(D|_{V\times_S P^\vee}\) 无固定点，亦即 \(Z\cap V=\varnothing\)，
从而 \(V=X\setminus Z\)。命题得证。\footnote{原稿由延拓态射直接援引命题3.1，省去了核验。更明确地，\((S_2)\) 保证几何纤维关联点都是分支泛点；每个 \(V\) 的几何纤维关联点已在 \(X\setminus Z\) 上。若整个关联分支映入某超平面，它在这个开集中的非空部分也映入该超平面，与命题3.1矛盾。或者先由关联除子对 \(V\) 的相对性得到一个不必对 \(P^\vee\) 平坦的有效 Cartier 除子，再用已证的延拓唯一性确认它等于原除子的限制。}

\clearpage
\section*{译校补证：非诺特底下如何使用深度}
原稿的局部化论证不应把任意闭集 \(T\) 不加论证地当成有限表示闭子概形。
下面只下降有限个环元素及有限表示的模，不下降整个 \(T\)。

在 \(x\in T\) 的仿射邻域令 \(I\) 定义 \(T\)，\(s\) 是 \(x\) 的像。纤维的局部环
\(\OO_{X_s,x}\) 诺特。在包含 \(I\) 的该局部环素点处，若纤维模非零，其深度都有假设中的下界；
因此由素避及正则元降深度，能在 \(I\) 的纤维像中找到长度一（或二）的纤维模正则序列。
若 \(\FF_x=0\)，有限生成性使它在某邻域已为零，可直接处理。
若 \(\FF_x\ne0\)，Nakayama 引理保证该点纤维模非零。
清去在该点可逆的分母，可把序列取为 \(I\) 中元素 \(h_1\)（或 \(h_1,h_2\)）的像。

对这些有限数据作诺特逼近：有限表示代数、有限表示模、底上的平坦性及这几个元素
在足够后的仿射阶段下降，见 \href{https://stacks.math.columbia.edu/tag/00R3}{有限表示模的平坦性下降}。
从较早阶段相应纤维的局部环到现有纤维局部环是忠实平坦同态，故正则性在该点下降。
在诺特阶段，\href{https://stacks.math.columbia.edu/tag/00ME}{纤维注入判据}
把逐次乘 \(h_j\) 的纤维注入提升为局部注入，并使逐次商模在底上平坦。
因诺特阶段核有限生成且商模的平坦处为开，缩小到该点的邻域后可同时实现这些断言；
再拉回，商模的平坦性保证短正合列仍正合。
所以在当前底上，局部也可找到这样的模正则序列。

长度一时，\(h_1\) 在模上单射，且主开集 \(D(h_1)\) 在 \(X\setminus T\) 内，
从而限制映射是单射。
长度二时，模 \(M\) 在 \(D(h_1)\cup D(h_2)\) 上的截面由 \(M\) 唯一延拓：
若两个分式 \(a/h_1^m,\ b/h_2^n\) 在交集相等，清去额外分母后可令 \(h_2^n a=h_1^m b\)。
由于 \(h_2\) 在 \(M/h_1^mM\) 上仍为非零因子，\(a=h_1^m c\)；
再由 \(h_1\) 的正则性得 \(b=h_2^n c\)。
因此先把原截面在这两个主开集上的限制延拓，再用 \(h_1\) 在模上单射，
可知该延拓在整个 \(X\setminus T\) 与原截面相同（均在所取的仿射邻域内）。
这就证明了所需的层同构，不需要把 \(X\setminus T\) 假定为拟紧。
\end{document}
